\section{Conclusion}
\label{sec:conclusion}

\draftnote{DRAFT placeholder. Expand once Sections~\ref{sec:residual-risk}
and~\ref{sec:phase4} are filled in -- the current text only restates
the introduction's two findings and should not be treated as final.}

Naive, group-agnostic evaluation of a learned mobility-privacy
mechanism's fairness is unreliable in (at least) two distinct ways we
have shown to be real and measurable, not merely plausible concerns:
the mechanism itself protects users unevenly, worst for the
lowest-entropy, most-identifiable decile (Section~\ref{sec:disparate-risk});
and a previously reported adaptive-adversary result that appeared to
show naive evaluation \emph{additionally} understates risk
asymmetrically does not survive decomposition into its actual causes
(Appendix~\ref{sec:adaptive-appendix}). Together these point toward a
single, narrower, and more defensible claim than ``this mechanism is
unfair'' or ``naive evaluation is meaningless'': individual-fairness
claims about a learned mobility-privacy mechanism must be measured at
the level of who is actually protected, not an aggregate or a naive
per-group average, and the adversary used to measure that protection
must itself be checked for confounds -- a smaller candidate pool or
side information an adversary would not realistically have -- before
its result is attributed to genuine sophistication.
